% ==============================================================================
% preamble.tex — 自由意思とフレーム慣性 (日本語版設定)
% ==============================================================================

\documentclass[11pt, a4paper, dvipdfmx]{bxjsarticle}

% ── 文字コード ───────────────────────────────────────────────────────────────
\usepackage[utf8]{inputenc}

% ── 页面布局 ──────────────────────────────────────────────────────────────
\geometry{margin=1in, headheight=14pt}
\usepackage{parskip}

% ── 数学符号 ──────────────────────────────────────────────────────────────
\usepackage{amsmath, amssymb, amsthm}

% ── 颜色与框 ──────────────────────────────────────────────────────────────
\usepackage[dvipsnames]{xcolor}
\usepackage[most]{tcolorbox}

% ── 超链接 ───────────────────────────────────────────────────────────────
\usepackage[dvipdfmx, colorlinks=true, linkcolor=RoyalBlue, urlcolor=RoyalBlue]{hyperref}

% ── 列表处理 ──────────────────────────────────────────────────────────────
\usepackage{enumitem}

% ── 页眉页脚 ──────────────────────────────────────────────────────────────
\usepackage{fancyhdr}
\pagestyle{fancy}
\fancyhf{}
\rhead{\textit{自由意思とフレーム慣性}}
\lhead{\textit{\nouppercase{\leftmark}}}
\cfoot{\thepage}

% ==============================================================================
% 定理環境 (日本語)
% ==============================================================================

\theoremstyle{definition}
\newtheorem{definition}{定義}[section]
\newtheorem{axiom}[definition]{公理}

\theoremstyle{plain}
\newtheorem{theorem}[definition]{定理}
\newtheorem{proposition}[definition]{命題}
\newtheorem{corollary}[definition]{系}
\newtheorem{lemma}[definition]{補題}

\theoremstyle{remark}
\newtheorem{remark}[definition]{注}

% ==============================================================================
% カスタム tcolorbox 環境
% ==============================================================================

\newtcolorbox{examplebox}[1][]{
  colback   = green!5!white,
  colframe  = green!60!black,
  fonttitle = \bfseries,
  title     = {例},
  #1
}

\newtcolorbox{discussion}[1][]{
  colback   = blue!3!white,
  colframe  = RoyalBlue!70!black,
  fonttitle = \bfseries,
  title     = {論考 / 解決},
  #1
}

\newtcolorbox{insightbox}[1][]{
  colback   = orange!5!white,
  colframe  = orange!80!black,
  fonttitle = \bfseries,
  title     = {重要視点},
  #1
}

\newtcolorbox{methodbox}[1][]{
  colback   = violet!5!white,
  colframe  = violet!70!black,
  fonttitle = \bfseries,
  title     = {習慣化メソッド},
  #1
}
